\section{Data and Simulations}\label{DatSim}

In the present analysis we have used the Planck data\footnote{Planck space telescope, operated by the ESA, had two instruments, the Low Frequency Instrument (LFI) and the High Frequency Instrument (HFI), with the goal of measuring the total intensity and polarization of CMB photons in a wide frequency range from 30 to 857 GHz (30 to 353 for polarization). This range was 
covered by 9 different frequencies, with the three lowest frequencies measured by the LFI (30, 44 and 70 GHz) and the rest by the HFI (100, 143, 217, 353, 545 and 857 GHz).} from the latest two releases, the Planck 2018 full-mission release (PR3) \cite{Planck:2018nkj}, and the fourth release (PR4), which has been processed by the NPIPE pipeline \cite{Planck:2020olo}. Both pipelines use the LFI and HFI raw and uncalibrated data
to generate frequency maps in the HEALPix \cite{Gorski:2004by} format, 9 in the case of temperature and 7 for Q and U Stokes parameters as the two highest frequency channels were not sensitive to polarization. These frequency maps are propagated through the component separation methods \cite{Planck:2018yye}, Sevem \cite{2012MNRAS.420.2162F}, SMICA \cite{2008arXiv0803.1814C}, NILC \cite{10.1111/j.1365-2966.2011.19770.x} and Commander \cite{2008ApJ...676...10E}, which generate each a clean CMB map. For PR3 we have used the data from the four separation algorithms. For PR4, clean CMB maps have been provided only by Sevem and Commander. We use both maps for the analysis in intensity and only Sevem products for polarization.\footnote{Although Commander polarization products are also available for PR4, 
we have not included them in this work, since the associated simulations were not well suited for our analysis (see Section \ref{PR4_E_Results}).} All the data and simulations are available at NERSC\footnote{National Energy Research Scientific Computing Center (NERSC), \url{https://www.nersc.gov/}, is a primary scientific computing facility operated by Lawrence Berkeley National Laboratory, located in California. It provides high-performance computing and storage facilities where Planck latest data and simulations can be found.} at full resolution\footnote{Data is also available at Planck Legacy Archive (PLA), \url{https:/pla.esac.esa.int/}.}  (N$_{\rm side} = 2048$), and have been downgraded (to N$_{\rm side}$=64) in harmonic space, according to, 
%
\begin{equation}\label{smoothing}
    a_{\ell m}^{out} = \frac{b_{\ell}^{out}p_{\ell}^{out}}{b_{\ell}^{in}p_{\ell}^{in}}a_{\ell m}^{in}
\end{equation}
%
where the superscripts $in$ and $out$ refer to the HEALPix resolution of the input and output maps, respectively, the $a_{\ell m}$'s are the harmonic coefficients of the considered maps, the $p_{\ell}$'s give the corresponding pixel window functions and the $b_{\ell}$'s are the Gaussian smoothing functions with a FWHM of 5$'$ for the original cleaned CMB maps and of 160$'$ for the downgraded maps. The output $a_{\ell m}$'s are then converted to a pixel map considering a maximum multipole $\ell_{max}$ = 3N$_{\rm side}-1$. For this procedure, we have used the Healpy package\footnote{\url{https://healpy.readthedocs.io/en/latest/}} \cite{Zonca2019}, which is basically a Python wrapper for HEALPix. Sevem maps for both
Planck releases
are shown in Figure \ref{fig:Maps}. Differences in polarization between both releases are apparent, mainly due to the reduction of large-scale systematics in the PR4 data. Conversely, intensity maps are very consistent.

All the details of the PR4 processing and the differences 
with respect to PR3
are given in \cite{Planck:2020olo}.
Most notably, while previous Planck pipelines process LFI and HFI in an independent way, PR4 was designed to process them simultaneously which produces a significant improvement in the calibration of both instruments. In addition, PR4 reduces the noise by including the data from the repointing maneuvers, which is around 8\% of the total mission. These differences, together with other added features and modifications,
have the net effect of reducing 
especially 
the polarization systematics at large scales 
as well as the level of noise. However, the signal at the smallest multipoles is also reduced due to the resulting transfer function (see section \ref{subsection:3.4} and Figure \ref{TF}).

\begin{figure}[t!]
\centering
  \includegraphics[scale = 0.43]{Figuras/Maps_PR3_PR4.png}
  \caption{Sevem component-separated CMB maps at 160$'$ resolution. The first column shows T, Q, U, and E-mode maps for the Planck 2018 pipeline (PR3), while the second column shows the same maps for NPIPE (PR4). E-mode maps are derived from the full-sky Stokes Q and U maps. The grey lines indicate either the temperature or the
  polarization Planck component separation confidence mask  \cite{Planck:2018yye}.}
  \label{fig:Maps}
\end{figure}

In order to test the null hypothesis, we have used as the reference point the FFP10 PR3 and NPIPE Monte Carlo (MC) simulations, provided by the Planck collaboration along with the data maps. This set of simulated maps includes realizations of the CMB signal, the instrumental noise and the systematics, 
which try to capture the characteristics of the 
full data processing
such as the scanning strategy, the detector responses, or the calibration errors introduced during data reduction. For both PR3 and PR4, 
a number of realizations of the CMB sky are generated according to a fiducial CMB power spectrum based on the cosmology described in Table 6 of \cite{Planck:2020olo}. These sky realizations include lensing, Rayleigh scattering, and Doppler boosting. Additionaly, they are convolved by the frequency-specific beam. 
In contrast, noise simulations have been obtained following a more complex approach (see section 4 of \cite{Planck:2019evm} and section 5 of \cite{Planck:2020olo}), processing them using the same algorithm as for the real data via an \emph{end-to-end} (E2E) pipeline, including instrumental noise and systematics. For PR3, after propagating E2E simulations through component-separation pipelines, the final product is a set of 999 CMB and 300 noise simulations for each of the four methods. On the other hand, for the PR4 pipeline, 600 CMB and noise realizations are available for Sevem,
while 100 simulations are provided for intensity in the case of Commander.

Regarding masks, we use the PR3 intensity and polarization component separation confidence masks downgraded to $N_{side} = 64$, which retains the 71.3\% and 72.4\% of the sky, respectively.
Following \cite{Planck:2019evm}, to generate lower resolution binary masks, we first smoothed them according to Eq. \ref{smoothing} and then set a threshold of 0.95. All the pixels with a value below the threshold are set to 0, while the remaining ones are set to 1. Masks at $N_{side} = 2048$ are available in the Planck Legacy Archive\footnote{\url{https://pla.esac.esa.int/}} (PLA). The intensity mask is used
without modification for the local-variance analysis in the temperature maps. On the other hand, the polarization mask is used as the starting point for the inpainting of the Q and U maps and to generate a set of extended customised masks for the analysis of the E-mode map, as explained in section \ref{subsection:3.2}. 